\documentclass[a4paper]{article}

% Basic packages
% Español (México) con XeLaTeX: fontspec para fuentes Unicode, polyglossia
% para la división silábica y los nombres del documento en la variante mexicana.
\usepackage{fontspec}
\usepackage{polyglossia}
\setdefaultlanguage[variant=mexican]{spanish}
% Los nombres chinos van como «pinyin (original)»: xeCJK da a los caracteres
% Han su propia fuente; sin ella XeLaTeX los omite con un aviso.
% FandolSong no cubre algunos caracteres de nombres (U+73FA); WenQuanYi Zen Hei sí.
\usepackage[AutoFallBack=true]{xeCJK}
\setCJKmainfont{FandolSong-Regular.otf}
\setCJKfallbackfamilyfont{\CJKrmdefault}{WenQuanYi Zen Hei}
% polyglossia no traduce los nombres de listings.
\AtBeginDocument{\providecommand{\lstlistingname}{}\renewcommand{\lstlistingname}{Listado}%
  \providecommand{\lstlistlistingname}{}\renewcommand{\lstlistlistingname}{Índice de listados}}
\usepackage{amsmath, amssymb}
\usepackage{graphicx}
\usepackage[margin=2.5cm]{geometry}
\usepackage[most]{tcolorbox}
\usepackage{etoolbox}
\usepackage{listings}
\usepackage{hyperref}
\usepackage{booktabs}
\usepackage{subcaption}
\usepackage{float}
\usepackage{tikz}
\IfFileExists{pgfplots.sty}{
    \usepackage{pgfplots}
    \pgfplotsset{compat=1.18}
}{}

% Highlight boxes
\newtcolorbox{knowledgebox}[1]{
    enhanced, colback=blue!5!white, colframe=blue!75!black, colbacktitle=blue!75!black,
    coltitle=white, fonttitle=\bfseries, title={#1}, attach boxed title to top left={yshift=-2mm, xshift=2mm},
    boxrule=1pt, sharp corners
}

\newtcolorbox{importantbox}[1]{
    enhanced, colback=yellow!10!white, colframe=yellow!80!black, colbacktitle=yellow!80!black,
    coltitle=black, fonttitle=\bfseries, title={#1}, sharp corners
}

\newtcolorbox{warningbox}[1]{
    enhanced, colback=red!5!white, colframe=red!75!black, colbacktitle=red!75!black,
    coltitle=white, fonttitle=\bfseries, title={#1}, sharp corners
}

% Listings
\lstset{
    language=Python,
    basicstyle=\ttfamily\small,
    keywordstyle=\color{blue},
    stringstyle=\color{red!60!black},
    commentstyle=\color{green!60!black},
    breaklines=true,
    frame=single,
    numbers=left,
    numberstyle=\tiny\color{gray},
    captionpos=b,
    extendedchars=true
}

% Front-page metadata
\newcommand{\notetitle}{AI for Biology: Diseño de proteínas impulsado por IA generativa}
\newcommand{\noteauthors}{Elaboradas a partir de la clase de Ava Amini}
\newcommand{\notedate}{Primavera de 2025}
\newcommand{\videochannel}{Microsoft / Ava Amini (Guest Lecture)}
\newcommand{\videopublishdate}{2025-05-05}
\newcommand{\videoduration}{57 minutos}
\newcommand{\videourl}{https://www.youtube.com/watch?v=SSzSOeGP87I}
\newcommand{\videocoverpath}{cover.jpg}
\newcommand{\repourl}{https://github.com/hqhq1025/ai-course-notes}

% Footer with repo info
\usepackage{fancyhdr}
\pagestyle{fancy}
\fancyhf{}
\fancyfoot[C]{\thepage}
\fancyfoot[R]{\footnotesize\color{gray}\href{\repourl}{AI Course Notes}}
\renewcommand{\headrulewidth}{0pt}
\renewcommand{\footrulewidth}{0pt}

\begin{document}

\begin{titlepage}
\centering
{\Large Notas del curso\par}
\vspace{1.2cm}
{\huge\bfseries \notetitle\par}
\vspace{0.8cm}
{\large \noteauthors\par}
\vspace{0.3cm}
{\large \notedate\par}
\vspace{1.2cm}

\ifdefempty{\videocoverpath}{
{\small Al generar las notas, indique la ruta local de la portada del video.\par}
}{
\includegraphics[width=0.82\textwidth,height=0.45\textheight,keepaspectratio]{\videocoverpath}\par
}

\vfill
\begin{tcolorbox}[width=0.9\textwidth, colback=black!2!white, colframe=black!60, sharp corners]
\textbf{Autor o canal del video}: \videochannel\par
\textbf{Fecha de publicación}: \videopublishdate\par
\textbf{Duración del video}: \videoduration\par
\ifdefempty{\videourl}{}{
\textbf{Enlace del video}: \href{\videourl}{\nolinkurl{\videourl}}\par
}\par
\textbf{Página del proyecto}: \href{\repourl}{\nolinkurl{\repourl}}\par
\end{tcolorbox}
\end{titlepage}

\tableofcontents
\newpage

%% --- Inicio del contenido --- %%

\section{Introducción: la convergencia de la IA y la biología}

Esta clase es la última del semestre de primavera de 2025 de MIT 6.S191 (Introduction to Deep Learning) y la imparte, como invitada especial, Ava Amini, investigadora de Microsoft Research (MSR). Amini tiene una formación que abarca tanto la biología de laboratorio húmedo como las ciencias computacionales: durante la licenciatura y el doctorado profundizó en experimentos celulares e investigación con modelos animales, y más adelante combinó su interés por la biología con métodos computacionales, dedicándose a impulsar la frontera de las ciencias de la vida mediante técnicas de IA.

\begin{knowledgebox}{La misión de Microsoft Research (MSR)}
La misión central de MSR es "avanzar la frontera de la ciencia y la tecnología en beneficio de la humanidad" (advance the frontiers of science and technology to benefit humanity). Su investigación se desarrolla desde tres niveles de perspectiva: (1) investigación de frontera, fundamental y de orientación académica; (2) despliegue responsable y éticamente adecuado de la tecnología; (3) un impacto positivo en la salud humana y en el descubrimiento científico en las ciencias naturales. Dentro de MSR existe un equipo dedicado, Biomedical Machine Learning (BioML), del cual Amini es integrante central.
\end{knowledgebox}

Amini señala que la biología es, en esencia, un sistema extremadamente complejo que opera a escala nanométrica. Incluso una visualización científica de la membrana de una sola célula, ampliada un millón de veces, revela una riqueza estructural sorprendente: una enorme cantidad de moléculas de proteína que desempeñan en ella todo tipo de funciones. Esta complejidad representa a la vez un enorme desafío para la IA y una oportunidad sin precedentes.

En el curso 6.S191, los estudiantes ya han aprendido los dos grandes paradigmas del \textbf{modelado predictivo} (predictive modeling) y el \textbf{modelado generativo} (generative modeling). Cuando este marco se aplica al ámbito biológico, la pregunta central sigue siendo la predicción y la generación, pero la modalidad de los datos cambia de manera fundamental:

\begin{itemize}
    \item Ya no se trata de texto en lenguaje natural, sino del \textbf{lenguaje de secuencias de las moléculas biológicas} individuales
    \item Ya no se trata de imágenes de rostros o de automóviles, sino de \textbf{representaciones celulares}, \textbf{imágenes de tejidos} y \textbf{datos de pacientes}
    \item El objetivo de predicción se convierte en: ¿cuál es la función de una molécula biológica? ¿Cómo responde una célula a un fármaco? ¿Cuál es el resultado clínico de un paciente?
    \item El objetivo de generación se convierte en: dada una función objetivo, ¿es posible diseñar una molécula biológica que la realice?
\end{itemize}

\begin{importantbox}{El paradigma bidireccional central de la IA para la biología}
\textbf{Predicción hacia adelante}: dada una molécula biológica $\rightarrow$ predecir su función (clasificación, regresión).\\
\textbf{Diseño inverso}: dada una función objetivo $\rightarrow$ generar una molécula biológica que la satisfaga (modelado generativo).\\
Estas dos direcciones se complementan y forman en conjunto el circuito completo mediante el cual la IA potencia a la biología. Lo fundamental es que tanto la predicción como el diseño de la IA deben, al final, conectarse con el mundo natural mediante \textbf{validación experimental}; esto es lo que distingue a la IA para la biología de los campos puramente computacionales.
\end{importantbox}

El contenido central de esta clase se centra en \textbf{la aplicación de la IA generativa al diseño de proteínas}, en particular en cómo trasladar el Diffusion Model (modelo de difusión), este potente marco de modelado generativo, del dominio de las imágenes al dominio de datos discretos de las secuencias de proteínas, y presenta el modelo \textbf{EvoDiff}, desarrollado por el equipo de Amini.

\subsection{Resumen de la sección}

La IA para la biología es una de las direcciones de aplicación de la IA con mayor potencial. Los dos pilares centrales de este campo son el modelado predictivo y el modelado generativo, y su particularidad radica en que necesariamente debe conectarse con el mundo natural mediante experimentación. Esta clase se centra en el uso de la IA generativa, en particular del Diffusion Model, para diseñar moléculas de proteína con una función objetivo.

\section{Fundamentos de las proteínas: secuencia, estructura y función}

\subsection{Las proteínas son los "actuadores" de la biología}

Las proteínas (protein) son una clase de macromoléculas biológicas extremadamente diversas, que pueden considerarse \textbf{los actuadores de la biología} (biology's actuators): impulsan y regulan casi todas las funciones dentro de la célula. El rango de funciones de las proteínas es sumamente amplio e incluye, entre otras:

\begin{itemize}
    \item \textbf{Enzimas} (enzyme): catalizan reacciones químicas específicas, incluida la edición a nivel del genoma (como las proteínas asociadas a CRISPR)
    \item \textbf{Proteínas estructurales}: constituyen el citoesqueleto, la matriz extracelular y otras estructuras
    \item \textbf{Proteínas de señalización}: transmiten señales entre células y dentro de ellas (como las proteínas que se unen al calcio)
    \item \textbf{Proteínas terapéuticas}: como las proteínas farmacológicas ya aprobadas clínicamente para tratar la leucemia
    \item \textbf{Proteínas de transporte}: transportan moléculas dentro y fuera de la célula
\end{itemize}

\begin{knowledgebox}{Las proteínas como fármacos terapéuticos}
Las proteínas no solo son ejecutoras naturales de funciones biológicas, también pueden \textbf{diseñarse mediante ingeniería} como fármacos terapéuticos. Por ejemplo, los medicamentos con anticuerpos (un tipo especial de proteína) se usan ampliamente en inmunoterapia contra el cáncer. Amini mostró en la clase un fármaco proteico ya aprobado clínicamente para tratar un tipo específico de leucemia, lo que demuestra que el valor de traslación clínica del diseño de proteínas ya se ha verificado.
\end{knowledgebox}

\subsection{Los tres niveles de representación de las proteínas}

La biología de las proteínas puede entenderse mediante tres niveles, que forman una relación jerárquica desde la codificación lineal hasta la estructura tridimensional y, finalmente, la función biológica:

\begin{table}[H]
\centering
\caption{Los tres niveles de representación de las proteínas}
\begin{tabular}{@{}lll@{}}
\toprule
\textbf{Nivel} & \textbf{Descripción} & \textbf{Características de los datos} \\
\midrule
Secuencia (Sequence) & Disposición lineal de aminoácidos & Secuencia de símbolos discretos, alfabeto de 20 aminoácidos \\
Estructura (Structure) & Conformación tridimensional en el espacio & Coordenadas continuas, información geométrica a nivel atómico \\
Función (Function) & Actividad biológica y efecto & Anotaciones diversas, difíciles de cuantificar \\
\bottomrule
\end{tabular}
\end{table}

\textbf{Nivel de secuencia}: cada proteína queda definida por una secuencia de aminoácidos (amino acid). En la naturaleza existen 20 aminoácidos estándar, cada uno representable con un código de una sola letra (por ejemplo, A para alanina, Alanine, y G para glicina, Glycine). Por ello, la secuencia de una proteína puede considerarse un "lenguaje" basado en un vocabulario de 20 caracteres, lo que ofrece una analogía natural para aplicar técnicas de PLN al procesamiento de datos de proteínas.

\textbf{Nivel de estructura}: la secuencia de aminoácidos se pliega (folding) mediante procesos fisicoquímicos complejos hasta adoptar una estructura tridimensional específica. La estructura determina cómo interactúa la proteína con otras moléculas y constituye la base física para la realización de su función.

\textbf{Nivel de función}: la función biológica de una proteína (como la catálisis, la unión o la transducción de señales) suele estar determinada por su estructura, y esta a su vez está codificada por la secuencia.

\begin{importantbox}{Ingeniería inversa del dogma central}
La investigación tradicional sobre proteínas sigue la vía directa de secuencia $\rightarrow$ estructura $\rightarrow$ función. Sin embargo, el diseño de proteínas impulsado por IA busca \textbf{invertir este paradigma}: partir de la \textbf{función} que deseamos obtener y deducir de manera inversa la \textbf{secuencia} que la realiza. Este es precisamente el reto central del diseño generativo de proteínas:
$$\text{Function} \rightarrow \text{Sequence} \rightarrow \text{Structure (验证)}$$
\end{importantbox}

\subsection{Herramientas de IA existentes para proteínas}

En el campo de la IA aplicada a proteínas ya existe una cantidad considerable de trabajo dirigido a cada uno de los eslabones del pipeline secuencia-estructura-función:

\begin{itemize}
    \item \textbf{Modelos de lenguaje de proteínas} (Protein Language Models): como la serie ESM, que aprenden representaciones de secuencias de proteínas y capturan información evolutiva y funcional
    \item \textbf{Modelos generativos de estructuras}: como RFdiffusion, capaces de diseñar en el espacio estructural nuevas geometrías que no existen en la naturaleza
    \item \textbf{Predicción de secuencia a estructura}: como \textbf{AlphaFold}, que dada una secuencia de aminoácidos predice su estructura tridimensional; este es uno de los avances más importantes de los últimos años en biología computacional
\end{itemize}

\begin{warningbox}{La brecha de escala entre los datos de estructura y los datos de secuencia}
En el campo de las proteínas existe una \textbf{asimetría de datos} considerable:
\begin{itemize}
    \item \textbf{Datos de secuencia}: cientos de millones, o incluso miles de millones, de secuencias de proteínas distintas (provenientes de la secuenciación de genomas)
    \item \textbf{Datos de estructura}: apenas unas 300,000 estructuras de proteínas resueltas experimentalmente (provenientes de la base de datos PDB)
\end{itemize}
Aunque AlphaFold puede predecir estructuras, se trata de una \textbf{predicción computacional} y no de una medición experimental. Las proteínas en las bases de datos de estructuras tienden a ser globulares (globular) y solubles, con abundancia de elementos estructurales de tipo $\alpha$-hélice, lo que introduce un sesgo sistemático en los modelos entrenados con datos de estructura.
\end{warningbox}

\subsection{Resumen de la sección}

Las proteínas son la macromolécula biológica central que conecta la codificación de secuencias con la función biológica. Comprender la relación jerárquica de tres niveles entre secuencia, estructura y función es la base para construir herramientas de IA de diseño de proteínas. Las herramientas existentes ya son capaces de abordar cada eslabón del pipeline por separado, pero lograr el diseño inverso de función a secuencia sigue siendo un reto abierto. En cuanto a los datos, la riqueza de los datos de secuencia supera con creces a la de los datos de estructura, y esta asimetría influye profundamente en la elección de las estrategias de diseño de los modelos.
\section{Repaso del Diffusion Model: de la imagen a la biología}

\subsection{Principio central del modelo de difusión}

Antes de profundizar en el diseño de proteínas, Amini primero repasó la idea central del Diffusion Model. El principio básico del modelo de difusión es \textbf{lograr la generación aprendiendo a recuperar los datos a partir del ruido}, e incluye dos procesos inversos entre sí:

\textbf{1. Proceso directo de adición de ruido (Forward Noising Process)}

Partiendo del espacio de datos original (como una imagen), se destruyen los datos de manera gradual mediante la adición iterativa de ruido:

$$x_0 \xrightarrow{\text{noise}} x_1 \xrightarrow{\text{noise}} x_2 \xrightarrow{\text{noise}} \cdots \xrightarrow{\text{noise}} x_T \sim \mathcal{N}(0, I)$$

Donde:
\begin{itemize}
    \item $x_0$: datos originales limpios
    \item $x_t$: datos después de agregar ruido en el paso $t$
    \item $x_T$: después de $T$ pasos, se aproxima a ruido puramente aleatorio
    \item Este proceso \textbf{no requiere entrenamiento}; simplemente se agrega ruido gaussiano de manera gradual
\end{itemize}

\textbf{2. Proceso inverso de eliminación de ruido (Reverse Denoising Process)}

Se entrena una red neuronal para aprender el proceso inverso, recuperando los datos de manera gradual a partir del ruido:

$$x_T \xrightarrow{\text{denoise}} x_{T-1} \xrightarrow{\text{denoise}} \cdots \xrightarrow{\text{denoise}} x_0$$

La tarea de entrenamiento puede formalizarse así: dado el dato con ruido $x_t$ en el paso de tiempo $t$, predecir una versión con menos ruido $x_{t-1}$ o, de manera equivalente, predecir el ruido residual $\epsilon$:

$$\mathcal{L} = \mathbb{E}_{x_0, \epsilon, t} \left[ \| \epsilon - \epsilon_\theta(x_t, t) \|^2 \right]$$

Donde:
\begin{itemize}
    \item $\epsilon$: el ruido realmente agregado
    \item $\epsilon_\theta$: el ruido predicho por la red neuronal
    \item $t$: el paso de tiempo muestreado aleatoriamente
\end{itemize}

\begin{importantbox}{¿Por qué el Diffusion Model es adecuado para el diseño de proteínas?}
El Diffusion Model ha mostrado tres ventajas revolucionarias en el campo de la generación de imágenes:
\begin{enumerate}
    \item \textbf{Calidad de generación}: el realismo y la fidelidad de las muestras son extremadamente altos
    \item \textbf{Diversidad}: puede cubrir una amplia región de la distribución de datos, sin sufrir mode collapse
    \item \textbf{Controlabilidad}: admite generación condicional, lo que permite especificar el objetivo de generación
\end{enumerate}
Estas características corresponden precisamente a las necesidades centrales del diseño de proteínas: las proteínas generadas deben ser biológicamente viables (calidad), cubrir funciones diversas (diversidad) y satisfacer especificaciones de diseño particulares (controlabilidad).
\end{importantbox}

\subsection{El reto de pasar del dominio continuo al discreto}

El Diffusion Model estándar procesa \textbf{datos continuos} (como los valores de los píxeles de una imagen), y el proceso directo se implementa agregando ruido gaussiano continuo. Sin embargo, la secuencia de proteínas es \textbf{datos discretos} —en cada posición hay uno de 20 aminoácidos posibles— y no se le puede simplemente «agregar ruido gaussiano» a un carácter.

Esto plantea un reto técnico central: \textbf{¿cómo definir el proceso de difusión sobre datos discretos?}

\subsection{Resumen de la sección}

El núcleo del Diffusion Model consiste en aprender el proceso inverso que va del ruido a los datos. Su éxito en el ámbito de las imágenes ofrece un paradigma sólido del cual tomar referencia para el diseño de proteínas, pero la naturaleza discreta de las secuencias de proteínas exige una transformación fundamental del marco estándar de difusión continua.

\section{Modelos de difusión discreta: adición y eliminación de ruido en secuencias de proteínas}

\subsection{Estrategias de adición de ruido para datos discretos}

Para datos de secuencias discretas (como el lenguaje natural o las secuencias de proteínas), Amini propuso dos estrategias para implementar la «adición de ruido»:

\textbf{Estrategia uno: Masking (enmascaramiento)}

Partiendo de una secuencia limpia, en cada paso se elige aleatoriamente una o varias posiciones y se sustituye el token correspondiente por la marca especial \texttt{[MASK]}:

\begin{table}[H]
\centering
\caption{Ilustración de la estrategia Masking (con el ejemplo en lenguaje natural «This class is fun»)}
\begin{tabular}{@{}cl@{}}
\toprule
\textbf{Paso} & \textbf{Estado de la secuencia} \\
\midrule
$t=0$ & This class is fun \\
$t=1$ & This {[}MASK{]} is fun \\
$t=2$ & {[}MASK{]} {[}MASK{]} is fun \\
$t=3$ & {[}MASK{]} {[}MASK{]} {[}MASK{]} fun \\
$t=4$ & {[}MASK{]} {[}MASK{]} {[}MASK{]} {[}MASK{]} \\
\bottomrule
\end{tabular}
\end{table}

\textbf{Estrategia dos: Mutation (mutación/sustitución)}

No se usa una marca especial; en cambio, en cada paso, con cierta probabilidad, se sustituyen aleatoriamente los tokens de algunas posiciones por otros tokens del vocabulario:

\begin{table}[H]
\centering
\caption{Ilustración de la estrategia Mutation}
\begin{tabular}{@{}cl@{}}
\toprule
\textbf{Paso} & \textbf{Estado de la secuencia} \\
\midrule
$t=0$ & This class is fun \\
$t=1$ & This \textit{apple} is fun \\
$t=2$ & \textit{Rock} \textit{apple} is fun \\
$t=3$ & \textit{Rock} \textit{apple} \textit{my} fun \\
$t=T$ & (secuencia de tokens completamente aleatoria) \\
\bottomrule
\end{tabular}
\end{table}

\begin{warningbox}{Mutation es más difícil que Masking}
La estrategia Mutation exige más del modelo de eliminación de ruido, porque este no puede identificar mediante una marca especial qué posiciones fueron modificadas: tiene que determinar por sí mismo qué tokens son «incorrectos» y predecir el valor correcto. Esto, en cierto sentido, se acerca más al proceso real de mutación evolutiva, pero también aumenta considerablemente la dificultad del entrenamiento. En EvoDiff se adoptaron y probaron ambas estrategias.
\end{warningbox}

\subsection{Formalización de la difusión discreta}

Para una secuencia de proteína $\mathbf{s} = (s_1, s_2, \ldots, s_L)$, donde $s_i \in \{A_1, A_2, \ldots, A_{20}\}$ representa uno de los 20 aminoácidos:

\textbf{Proceso hacia adelante}: en cada paso $t$, se elige aleatoriamente un subconjunto de posiciones $\mathcal{M}_t$ y se sustituyen los aminoácidos de esas posiciones por \texttt{[MASK]}:

$$s_i^{(t)} = \begin{cases} \texttt{[MASK]} & \text{if } i \in \mathcal{M}_1 \cup \mathcal{M}_2 \cup \cdots \cup \mathcal{M}_t \\ s_i^{(0)} & \text{otherwise} \end{cases}$$

\textbf{Proceso inverso}: se entrena una red neuronal $f_\theta$ que, en cada paso, restaura algunas de las posiciones enmascaradas a su aminoácido correcto:

$$f_\theta(\mathbf{s}^{(t)}, t) \rightarrow \mathbf{s}^{(t-1)}$$

Al generar, se parte de la secuencia totalmente enmascarada $\mathbf{s}^{(T)} = (\texttt{[MASK]}, \texttt{[MASK]}, \ldots, \texttt{[MASK]})$ y se elimina el ruido paso a paso hasta obtener una secuencia de proteína completa.

\subsection{Relación entre la difusión discreta y los modelos de lenguaje}

Amini señaló una conexión teórica profunda: \textbf{la difusión discreta con Masking es, en realidad, una generalización de varios esquemas de modelado de lenguaje}:

\begin{table}[H]
\centering
\caption{Comparación entre la difusión discreta y los esquemas clásicos de modelado de lenguaje}
\begin{tabular}{@{}lll@{}}
\toprule
\textbf{Método} & \textbf{Orden de decodificación} & \textbf{Predicción por paso} \\
\midrule
Next Token Prediction (LLM) & Fijo, de izquierda a derecha & Un solo token siguiente \\
Masked Language Model (BERT) & Toda la secuencia visible & Se rellenan todos los mask en un paso \\
Discrete Diffusion & Orden aleatorio, paso a paso & Se rellena parte de los mask en cada paso \\
\bottomrule
\end{tabular}
\end{table}

\begin{importantbox}{Capacidad de generalización de la difusión discreta}
El modelo de difusión discreta aprende la tarea de eliminación de ruido sobre \textbf{todos los órdenes de decodificación} y \textbf{todos los subconjuntos de masking} posibles. Esto hace que abarque, al mismo tiempo:
\begin{itemize}
    \item la decodificación «de izquierda a derecha» de los modelos autorregresivos (un orden de masking particular)
    \item la «eliminación de ruido en un paso» de Masked Language Model (una configuración particular del número de pasos)
\end{itemize}
Por lo tanto, el modelo de difusión discreta es un marco más general, con mayor flexibilidad y capacidad expresiva.
\end{importantbox}

\subsection{Resumen de la sección}

Trasladar el Diffusion Model del dominio continuo al dominio discreto requiere redefinir el «ruido»: mediante Masking o Mutation se corrompen progresivamente los tokens de la secuencia. El marco de difusión discreta unifica matemáticamente los modelos autorregresivos y Masked Language Model, y ofrece una herramienta poderosa y flexible para el modelado de secuencias de proteínas.


\section{EvoDiff: un modelo de difusión de secuencias de proteínas a escala evolutiva}

\subsection{Panorama del modelo y datos de entrenamiento}

EvoDiff es un \textbf{modelo generativo de proteínas basado en difusión discreta} desarrollado por el equipo de Amini en MSR, construido específicamente para el diseño de proteínas funcionales. El proyecto fue dirigido por Sarah Alamdari, con Amini y su colega Kevin Yang codirigiendo la orientación de la investigación.

\begin{importantbox}{Características centrales de EvoDiff}
\begin{itemize}
    \item \textbf{Tipo de modelo}: modelo generativo basado en difusión discreta (Discrete Diffusion)
    \item \textbf{Modalidad de datos}: secuencias de aminoácidos de proteínas (tokens discretos)
    \item \textbf{Datos de entrenamiento}: alrededor de 50 millones de secuencias de proteínas únicas, a escala evolutiva
    \item \textbf{Capacidad central}: generar nuevas secuencias de proteínas eliminando el ruido gradualmente a partir de un estado totalmente enmascarado
    \item \textbf{Capacidad extendida}: admite generación condicionada por función (motif scaffolding)
\end{itemize}
\end{importantbox}

«A escala evolutiva» significa que los datos de entrenamiento abarcan las secuencias de proteínas de distintos organismos a lo largo del árbol de la vida (tree of life), desde bacterias hasta humanos, e incluyen una representación de funciones biológicas extremadamente diversa. Esta amplia cobertura de datos permite que el modelo aprenda las regularidades generales de la distribución de las proteínas naturales.

\subsection{Dos modos de entrada de secuencia}

EvoDiff admite dos modos de entrada de secuencia, cada uno adecuado para un escenario de modelado distinto:

\textbf{Modo uno: secuencia única (Single Sequence)}

Se modela directamente una sola secuencia de proteína mediante difusión discreta. Es adecuado para el diseño de proteínas desde cero (de novo design) sin condicionamiento.

\textbf{Modo dos: Multiple Sequence Alignment (MSA)}

Aprovecha la información de contexto evolutivo. Para una secuencia objetivo dada, mediante una búsqueda de similitud de secuencias se encuentra un conjunto de \textbf{secuencias homólogas evolutivas relacionadas pero no idénticas}, estas secuencias se alinean y se disponen en una matriz (MSA) y, después, se ejecuta la difusión discreta sobre toda la matriz de MSA.

\begin{knowledgebox}{¿Por qué el MSA aporta información adicional?}
Multiple Sequence Alignment es un concepto clásico de la biología computacional. Las secuencias de proteínas relacionadas evolutivamente, aunque presentan mutaciones en ciertas posiciones, suelen mantenerse muy conservadas en las posiciones clave para la función. El patrón de conservación (conservation pattern) a nivel de columna en el MSA codifica de manera implícita la información crucial de \textbf{qué posiciones son fundamentales para la función}. El modo MSA de EvoDiff, al modelar de forma conjunta varias secuencias relacionadas, puede extraer un conocimiento previo más rico a partir de las señales evolutivas y así guiar el proceso de generación.
\end{knowledgebox}

\subsection{Visualización del proceso de generación}

El proceso de generación de EvoDiff puede entenderse de manera intuitiva así:

\begin{enumerate}
    \item \textbf{Inicialización}: una secuencia completamente en \texttt{[MASK]}, con la longitud de la proteína objetivo
    \item \textbf{Eliminación iterativa del ruido}: en cada paso, el modelo predice y rellena los aminoácidos de una parte de las posiciones enmascaradas
    \item \textbf{Formación de la secuencia}: a medida que avanzan los pasos, se rellenan cada vez más posiciones y la secuencia va tomando forma gradualmente
    \item \textbf{Finalización}: se rellenan todas las posiciones y se obtiene una secuencia de proteína completa
\end{enumerate}

Amini mostró en la clase una visualización dinámica: del lado izquierdo, el proceso de generación paso a paso del modelo de secuencias de EvoDiff; del lado derecho, sincronizada, la estructura tridimensional prevista que corresponde a la secuencia actual. Vale la pena enfatizar que \textbf{el aprendizaje y la generación de EvoDiff ocurren por completo en el espacio de secuencias}, sin utilizar ninguna información estructural; la estructura del lado derecho solo se usa para la visualización final.

\subsection{Resumen de la sección}

EvoDiff utiliza un marco de difusión discreta, entrenado sobre alrededor de 50 millones de secuencias de proteínas a escala evolutiva, y admite los dos modos de entrada de secuencia única y MSA. El modelo opera por completo en el espacio de secuencias, generando nuevas proteínas al eliminar el ruido gradualmente a partir de un estado totalmente enmascarado. El modo MSA, al incorporar información de contexto evolutivo, aporta un conocimiento previo más rico al proceso de generación.
\section{Evaluación de EvoDiff: calidad, diversidad y validación funcional}

\subsection{Metodología de evaluación}

Evaluar un modelo de generación de proteínas es mucho más complejo que evaluar un modelo de generación de imágenes. Amini enfatiza que es necesario diseñar la estrategia de evaluación con cuidado desde \textbf{múltiples niveles}:

\begin{table}[H]
\centering
\caption{Marco de evaluación multinivel para modelos de generación de proteínas}
\begin{tabular}{@{}lll@{}}
\toprule
\textbf{Dimensión de evaluación} & \textbf{Pregunta central} & \textbf{Método de evaluación} \\
\midrule
Calidad individual & ¿Es la proteína individual biológicamente viable? & Predicción de estructura + self-consistency \\
Cobertura de la distribución & ¿Cubren las muestras generadas la distribución de datos? & Visualización y comparación de la distribución en el espacio de características \\
Realización funcional & ¿Puede la proteína generada cumplir la función buscada? & Validación experimental en laboratorio húmedo \\
\bottomrule
\end{tabular}
\end{table}

\begin{warningbox}{No basta con ver la "exactitud"}
La evaluación de un modelo de generación de proteínas \textbf{no puede reducirse a un solo indicador de exactitud}. Un modelo puede generar proteínas individuales de alta calidad, pero si todas las proteínas generadas se concentran en una región estrecha del espacio funcional, su valor como herramienta de diseño se ve muy disminuido. Por eso, la \textbf{diversidad a nivel de distribución} y la \textbf{capacidad de validación a nivel funcional} son tan importantes como la calidad individual.
\end{warningbox}

\subsection{Evaluación de la calidad individual: Self-Consistency}

EvoDiff adopta un método ingenioso de evaluación por self-consistency:

\begin{enumerate}
    \item Se genera con EvoDiff una nueva secuencia de proteína $\mathbf{s}_{\text{gen}}$
    \item Se predice con una herramienta de predicción de estructura (como AlphaFold) la estructura tridimensional $\mathcal{S}$ de $\mathbf{s}_{\text{gen}}$
    \item A partir de la estructura $\mathcal{S}$, se infiere una secuencia $\mathbf{s}_{\text{inv}}$ capaz de plegarse en esa estructura
    \item Se calcula la consistencia entre $\mathbf{s}_{\text{gen}}$ y $\mathbf{s}_{\text{inv}}$
\end{enumerate}

Si la consistencia es alta, indica que la secuencia generada puede plegarse de manera estable en una estructura tridimensional determinada, es decir, es \textbf{estructuralmente viable} (structurally realistic and consistent). Los resultados muestran que las proteínas generadas por EvoDiff tienen un buen desempeño en consistencia estructural.

\subsection{Evaluación de la cobertura de la distribución: comparación con proteínas naturales}

Para evaluar la diversidad de las muestras generadas por EvoDiff, el equipo realizó un análisis a nivel de distribución:

\begin{enumerate}
    \item Se toma el conjunto de prueba de secuencias de proteínas naturales y se extraen vectores de características con un modelo preentrenado
    \item Los vectores de características se proyectan a un espacio 2D mediante un método de reducción de dimensionalidad (como t-SNE o UMAP), obteniendo la distribución de las proteínas naturales
    \item Se realiza la misma operación con una gran cantidad de proteínas generadas por EvoDiff (miles de secuencias)
    \item Se superponen ambas distribuciones para evaluar el grado de cobertura
\end{enumerate}

\begin{importantbox}{Ventaja de EvoDiff en la cobertura de la distribución}
La comparación con varios métodos de referencia muestra una ventaja notable de EvoDiff:
\begin{itemize}
    \item \textbf{frente a modelos de Next Token Prediction}: ambos alcanzan una cobertura aproximadamente equivalente, pero el modelo de next token es ligeramente superior
    \item \textbf{frente a modelos de generación en un solo paso con Masked}: la cobertura de EvoDiff es claramente más amplia, lo que indica que el denoising iterativo es superior a la generación en un solo paso
    \item \textbf{frente a métodos puramente estructurales} (como RFdiffusion, entre otros): los métodos estructurales muestran un \textbf{sesgo grave}, con una gran cantidad de muestras concentradas en una región estrecha del espacio funcional
\end{itemize}
\end{importantbox}

El sesgo grave de los métodos puramente estructurales tiene su origen en la distribución sesgada de los datos estructurales: entre las aproximadamente 300,000 estructuras de proteínas resueltas experimentalmente, las proteínas globulares y las ricas en $\alpha$-helix están sobrerrepresentadas. Los modelos entrenados con estos datos naturalmente sobremuestrean (oversample) este tipo de proteínas. EvoDiff, en cambio, se entrena con datos de secuencia, cuyo volumen es mayor y cuya cobertura es más amplia, por lo que captura mejor la diversidad del espacio funcional de las proteínas naturales.

\subsection{Validación experimental: de lo computacional al laboratorio}

La evaluación de referencia definitiva es la \textbf{validación experimental en laboratorio húmedo}. El equipo tomó proteínas candidatas de entre los resultados generados por EvoDiff y realizó los siguientes experimentos:

\textbf{Validación de la estabilidad estructural}:
\begin{itemize}
    \item Se sometieron secuencias de proteínas diseñadas por completo de manera artificial a \textbf{expresión biológica} (síntesis de la proteína dentro de la célula)
    \item Se midió la estabilidad estructural de la proteína
    \item Resultado: se expresaron y validaron con éxito 4 diseños de proteínas completamente nuevas y estructuralmente estables
\end{itemize}

\begin{knowledgebox}{El ciclo completo, del diseño computacional a la validación experimental}
El valor final del diseño de proteínas con IA debe demostrarse mediante validación experimental. Esto implica un proceso complejo: (1) el modelo computacional genera secuencias candidatas; (2) mediante síntesis génica se obtiene el ADN correspondiente; (3) la proteína se expresa en un sistema celular; (4) se purifica la proteína; (5) se valida su estructura y función con métodos fisicoquímicos (como dicroísmo circular, cristalografía de rayos X, etc.). El equipo de EvoDiff completó todo este ciclo, demostrando la utilidad práctica del modelo.
\end{knowledgebox}

\subsection{Resumen de la sección}

La evaluación de EvoDiff abarcó tres niveles: calidad individual, diversidad de la distribución y validación experimental. En cuanto a la cobertura de la distribución, EvoDiff muestra una ventaja clara frente a los métodos puramente estructurales, lo cual se debe a la ventaja intrínseca de los datos de secuencia en escala y diversidad. Lo más importante es que las proteínas generadas por EvoDiff se expresaron y se validaron estructuralmente con éxito en el laboratorio, lo que demuestra la utilidad práctica del modelo.

\section{Generación condicional funcional: Motif Scaffolding e indicaciones biológicas}

\subsection{De la generación incondicional al control funcional}

El verdadero poder del marco EvoDiff radica en la \textbf{generación condicional funcional} (conditional generation): no solo puede generar proteínas desde cero, sino también guiar el proceso de generación de acuerdo con restricciones funcionales específicas. Esto se debe a la flexibilidad natural del marco de difusión discreta.

La idea central es \textbf{Motif Scaffolding} (diseño de un andamiaje para un motivo funcional):

\begin{importantbox}{Concepto central de Motif Scaffolding}
La función clave de muchas proteínas está determinada por una \textbf{subregión local} de la secuencia (llamada motivo o motivo funcional). Por ejemplo, en una proteína de unión a iones de calcio, solo unos cuantos residuos de aminoácidos específicos participan directamente en la unión física con el ion de calcio.

\textbf{El objetivo de Motif Scaffolding}: dado un motivo funcional conocido, diseñar una secuencia de proteína completamente nueva que «sostenga» (scaffold) ese motivo, de modo que pueda cumplir su función en la posición tridimensional correcta.

Esto es similar a la tarea de \textbf{inpainting} (relleno) en el procesamiento de lenguaje natural: dada una parte del contenido de la secuencia, se le pide al modelo que complete el resto.
\end{importantbox}

\subsection{Inpainting como indicación funcional}

La capacidad de Motif Scaffolding de EvoDiff proviene directamente de su objetivo de entrenamiento. Dado que el modelo aprendió a eliminar el ruido de \textbf{todos los patrones posibles de masking}, adquiere de manera natural las siguientes capacidades:

\begin{enumerate}
    \item Recibir una secuencia \textbf{parcialmente conocida}: la parte conocida es el motivo funcional, y el resto de las posiciones son \texttt{[MASK]}
    \item Eliminar el ruido de forma gradual solo en las posiciones enmascaradas, manteniendo intacta la parte del motivo
    \item Generar finalmente una secuencia de proteína que contiene el motivo original pero cuyo resto es completamente nuevo
\end{enumerate}

Amini compara esta capacidad con una \textbf{«indicación biológica»} (biological prompting): el motivo funcional es el «prompt» que se le da al modelo, y este genera, a partir de ese prompt, una proteína completa que satisface la restricción funcional.

\subsection{Estudio de caso: diseño de una proteína de unión a iones de calcio}

Amini presentó en detalle un caso concreto de diseño funcional:

\textbf{Contexto}: existe de forma natural una proteína cuya función es unirse a iones de calcio (Ca$^{2+}$) y transmitir señales celulares. En la secuencia completa de esta proteína, hay una subsecuencia específica (el motivo marcado en verde) que es directamente responsable de la unión física y la interacción con el ion de calcio.

\textbf{Proceso de diseño}:
\begin{enumerate}
    \item Extraer el fragmento de secuencia del motivo de unión a calcio como «prompt»
    \item Colocar el motivo en la posición correspondiente de la secuencia objetivo, y fijar el resto de las posiciones como \texttt{[MASK]}
    \item Ejecutar el flujo de generación condicional de EvoDiff, con el que el modelo rellena de forma gradual las posiciones enmascaradas
    \item Obtener un diseño de proteína que contiene el motivo original de unión a calcio, pero cuya secuencia global es completamente nueva
\end{enumerate}

\textbf{Validación experimental}:
\begin{itemize}
    \item Se sometió la proteína diseñada por EvoDiff (curva azul) y la proteína natural (curva gris) a un experimento de unión a iones de calcio
    \item Los resultados muestran que la proteína diseñada por IA \textbf{exhibe con éxito capacidad de unión a iones de calcio}
    \item Esto significa que EvoDiff no solo puede generar proteínas viables estructuralmente, sino también proteínas \textbf{viables funcionalmente}
\end{itemize}

\begin{warningbox}{La validación funcional aún se encuentra en una etapa temprana}
Amini reconoció con franqueza que, si bien los experimentos actuales de validación funcional son alentadores, siguen siendo una «prueba de concepto» (proof of concept) preliminar. Del éxito de un solo caso a una capacidad sistemática de diseño funcional aún queda una cantidad considerable de trabajo por delante. El ciclo de validación experimental del diseño de proteínas es largo y costoso, y detrás de cada caso exitoso puede haber múltiples intentos fallidos.
\end{warningbox}

\subsection{Resumen de la sección}

El marco de difusión discreta de EvoDiff admite de manera natural la generación condicional funcional (Motif Scaffolding): al usar un motivo funcional conocido como «indicación biológica», guía al modelo para generar una proteína completamente nueva que contenga dicho motivo. El caso de diseño de la proteína de unión a iones de calcio muestra que la proteína generada por EvoDiff no solo es viable estructuralmente, sino que además logra la función deseada. Esto marca un primer paso importante hacia el diseño de proteínas funcionales impulsado por IA.
\section{De las moléculas a las células: la investigación en biología con IA a través de escalas}

\subsection{Un panorama más amplio que el del diseño de proteínas}

Amini enfatiza que el diseño de proteínas es solo un nivel dentro de este panorama más amplio de la IA en biología. La biología es un \textbf{sistema profundamente jerarquizado}:

\begin{table}[H]
\centering
\caption{La estructura jerárquica de la biología y los retos de IA correspondientes}
\begin{tabular}{@{}llll@{}}
\toprule
\textbf{Nivel} & \textbf{Objeto} & \textbf{Tipo de datos} & \textbf{Tarea de IA} \\
\midrule
Molecular & Proteínas, ADN & Secuencia, estructura & Diseño molecular, predicción de función \\
Celular & Estado de la célula individual & Transcriptoma, proteoma & Caracterización del estado celular, predicción de perturbaciones \\
Tisular & Cortes de tejido & Imágenes patológicas & Detección, clasificación y pronóstico de enfermedades \\
Individual & Datos del paciente & Multiómica, registros clínicos & Medicina de precisión, predicción de eficacia terapéutica \\
\bottomrule
\end{tabular}
\end{table}

La investigación del equipo MSR BioML abarca todos estos niveles, con el objetivo de desarrollar métodos de IA capaces de razonar en distintas escalas y, finalmente, integrarlos para impulsar decisiones con impacto clínico.

\subsection{La dirección de fusión entre secuencia y estructura}

En el campo del diseño de proteínas, el equipo está explorando activamente el \textbf{modelado conjunto} de la información de secuencia y de estructura. Amini señala que secuencia y estructura son complementarias:

\begin{itemize}
    \item La \textbf{estructura} proporciona un control geométrico preciso de alta resolución
    \item La \textbf{secuencia} proporciona una cobertura más amplia de la diversidad biológica
\end{itemize}

Entre las vías para combinar ambas están: mapear la secuencia y la estructura a un espacio de representación común, y entrenar conjuntamente modelos generativos de secuencia-estructura, entre otras.

\subsection{Multimodalidad texto-proteína}

Otra dirección de vanguardia es la \textbf{alineación multimodal} entre proteínas y lenguaje natural:

\begin{itemize}
    \item Dada una secuencia/estructura de proteína $\rightarrow$ generar un texto de descripción funcional
    \item Dado un texto de descripción funcional $\rightarrow$ diseñar una proteína que satisfaga esa descripción
\end{itemize}

Esto avanza hacia la visión que Amini describió al inicio de la conferencia: usar un prompt en lenguaje natural (como «diseña una proteína capaz de dirigirse a células de cáncer de mama y unirse a ellas») para impulsar el diseño de proteínas.

\subsection{IA a escala celular: entender los estados de enfermedad}

A nivel celular, el equipo está desarrollando métodos de IA capaces de definir y caracterizar los \textbf{estados celulares} (cellular states). El objetivo es entender cómo la enfermedad (en particular el cáncer) modifica el estado celular, y usar esa información para orientar la selección de estrategias terapéuticas.

\subsection{Patología digital: análisis de imágenes de tejido impulsado por IA}

A nivel tisular, el equipo de MSR ha logrado avances decisivos en el campo de la \textbf{patología digital} (digital pathology): la aplicación de métodos de visión por computadora al análisis de imágenes histopatológicas de tejido, para lograr la detección, clasificación y pronóstico automáticos del cáncer.

\begin{knowledgebox}{La colaboración estratégica entre MSR y el Broad Institute}
El equipo MSR BioML ha establecido una colaboración de investigación estratégica y de largo plazo con el Broad Institute del MIT y Harvard. El Broad Institute es una institución líder a nivel mundial en investigación de medicina de precisión y genómica. El objetivo de esta colaboración es articular estrechamente las capacidades de IA con recursos experimentales y clínicos de primer nivel, para lograr un ciclo completo que va de los datos moleculares del paciente a la predicción de IA, de ahí a la validación experimental y, finalmente, a la aplicación clínica. Este modelo de colaboración refleja el rasgo central del campo de la IA para la biología: la imprescindible complementariedad entre lo computacional y lo experimental.
\end{knowledgebox}

\subsection{Resumen de la sección}

La investigación en IA para la biología va mucho más allá del diseño de proteínas: es un proyecto sistemático que atraviesa múltiples niveles, desde el molecular hasta el celular, el tisular y el individual. El trabajo del equipo MSR BioML cubre desde el modelado conjunto de secuencia y estructura de proteínas hasta la patología digital, con la meta final de integrar métodos de IA multiescala para impulsar decisiones de medicina de precisión con impacto real en la práctica clínica.
\section{Preguntas y respuestas destacadas: perspectivas clave en las preguntas profundas}

\subsection{Origen del sesgo en los métodos estructurales}

\textbf{P}: ¿Por qué la distribución de generación de los métodos puramente estructurales presenta un sesgo severo?

\textbf{R}: La causa fundamental es el \textbf{sesgo de muestreo de los datos estructurales}. Las proteínas con estructura resuelta experimentalmente (alrededor de 300,000 en la base de datos PDB) tienden a presentar las siguientes características:
\begin{itemize}
    \item \textbf{proteínas globulares} (globular proteins): compactas, solubles en disoluciones acuosas, fáciles de manejar experimentalmente
    \item \textbf{ricas en $\alpha$-helix}: este tipo de elemento estructural está sobrerrepresentado entre las proteínas ya resueltas
\end{itemize}

El artículo de EvoDiff aporta datos detallados al respecto: las proteínas generadas por EvoDiff cubren mejor otros elementos estructurales además de la $\alpha$-helix (como la $\beta$-sheet, entre otros), mientras que los métodos puramente estructurales tienden a sobremuestrear proteínas ricas en $\alpha$-helix.

\subsection{Perspectivas de la fusión de secuencia y estructura}

\textbf{P}: ¿Podrá EvoDiff integrar información estructural en el futuro?

\textbf{R}: Sí, esta es una dirección de investigación activa en el campo del aprendizaje automático aplicado a proteínas. Las estrategias concretas de fusión incluyen:
\begin{itemize}
    \item mapear la secuencia y la estructura a un \textbf{espacio de representación común}
    \item construir un \textbf{modelo generativo conjunto de secuencia y estructura}
    \item usar la información estructural para \textbf{restringir y guiar} la generación de la secuencia
\end{itemize}

Las dos modalidades de datos son \textbf{complementarias}: los datos estructurales son escasos en volumen pero aportan restricciones geométricas precisas, mientras que los datos de secuencia son abundantes en volumen y cubren una amplia diversidad biológica. Fusionarlos promete reunir las ventajas de ambos.

\subsection{Resumen de la sección}

La sesión de preguntas y respuestas reveló dos perspectivas técnicas importantes: (1) el sesgo sistemático de los datos estructurales es la limitación fundamental de los métodos puramente estructurales; (2) la fusión de secuencia y estructura es el camino necesario para lograr capacidades de diseño de proteínas más sólidas.

\section{Síntesis y ampliación}

\subsection{Repaso de los puntos clave}

Esta clase partió de la intersección entre la IA y la biología para presentar en profundidad las aplicaciones de vanguardia de la IA generativa en el diseño de proteínas; el contenido central puede resumirse en los siguientes niveles:

\begin{enumerate}
    \item \textbf{La biología es uno de los campos de aplicación con mayor impacto para la IA}: las proteínas, como «ejecutoras» de la biología, y los avances en la capacidad de diseñarlas influirán profundamente en el desarrollo de fármacos, la ingeniería de enzimas, la biología sintética y otros campos

    \item \textbf{La jerarquía de tres niveles secuencia-estructura-función de las proteínas} es el marco básico para entender y manipular las proteínas; el diseño de proteínas impulsado por IA busca invertir esta jerarquía, diseñando la secuencia de manera inversa a partir de la función

    \item \textbf{El modelo de difusión (Diffusion Model) puede trasladarse del dominio continuo al dominio discreto}: mediante estrategias de Masking o Mutation se define el proceso de adición de ruido sobre datos discretos, con lo que se establece un marco completo de difusión discreta

    \item \textbf{EvoDiff es un modelo de generación de secuencias de proteínas basado en difusión discreta}: se entrenó con aproximadamente 50 millones de secuencias a escala evolutiva y admite tanto la generación incondicional como la generación condicionada por función (Motif Scaffolding)

    \item \textbf{La validación experimental es el criterio final}: las proteínas generadas por EvoDiff ya se han expresado con éxito en el laboratorio, y se ha verificado tanto su estabilidad estructural como su actividad funcional (unión de iones de calcio)

    \item \textbf{AI for Biology es una ingeniería de sistemas que abarca múltiples escalas}: desde la molécula hasta la célula, el tejido y el individuo, se requiere la integración y colaboración de métodos de IA multimodal
\end{enumerate}

\subsection{El significado teórico de los modelos de difusión discreta}

Desde la perspectiva de la teoría del aprendizaje profundo, esta clase reveló una perspectiva unificadora profunda: el marco de Masking Diffusion discreto unifica matemáticamente los modelos de lenguaje autorregresivos y el Masked Language Model, y constituye un paradigma de modelado generativo más general. Esta contribución teórica no solo tiene valor para el diseño de proteínas, sino que también resulta esclarecedora para tareas más amplias de generación de datos discretos (como el diseño de moléculas, la generación de código, etc.).

\subsection{Retos clave y direcciones futuras}

Aunque EvoDiff mostró resultados iniciales apasionantes, el campo todavía enfrenta retos importantes:

\begin{itemize}
    \item \textbf{Escalabilidad de la validación funcional}: la validación experimental actual todavía se limita al nivel de casos individuales; cómo escalarla al nivel de sistemas es un cuello de botella clave
    \item \textbf{Modelado conjunto de secuencia, estructura y función}: cómo fusionar de manera eficaz la información de distintas modalidades y, al mismo tiempo, manejar las enormes diferencias en la escala de los datos
    \item \textbf{Interfaces de diseño impulsadas por lenguaje natural}: lograr un sistema de extremo a extremo que «describa los requisitos funcionales en lenguaje natural $\rightarrow$ genere automáticamente la proteína que cumpla esos requisitos»
    \item \textbf{Optimización experimental en ciclo cerrado}: establecer un ciclo iterativo cerrado de diseño con IA $\rightarrow$ prueba experimental $\rightarrow$ optimización por retroalimentación (active learning)
    \item \textbf{Integración entre escalas}: modelado de extremo a extremo desde el diseño molecular hasta la respuesta celular y los efectos a nivel de tejido o paciente
\end{itemize}

\subsection{Lecturas adicionales}

\begin{itemize}
    \item \textbf{Artículo de EvoDiff}：Alamdari et al., ``Protein generation with evolutionary diffusion: sequence is all you need'', \textit{Nature Biotechnology}, 2023
    \item \textbf{AlphaFold}：Jumper et al., ``Highly accurate protein structure prediction with AlphaFold'', \textit{Nature}, 2021
    \item \textbf{RFdiffusion} (método de difusión estructural)：Watson et al., ``De novo design of protein structure and function with RFdiffusion'', \textit{Nature}, 2023
    \item \textbf{Modelo de lenguaje de proteínas ESM}：Lin et al., ``Evolutionary-scale prediction of atomic-level protein structure with a language model'', \textit{Science}, 2023
    \item \textbf{Revisión sobre el modelado conjunto de secuencia y estructura}：Wang et al., artículo de revisión del equipo de MSR (mencionado en la clase)
    \item \textbf{Página del curso MIT 6.S191}：\url{http://introtodeeplearning.com/}
    \item \textbf{Equipo de MSR BioML}：\url{https://www.microsoft.com/en-us/research/}
\end{itemize}

%% --- Fin del contenido --- %%

\end{document}
